% DRAFT — Han Liu: check every item below, replace anything that is not
% actually yours, delete the ones you did not do, then commit this file under
% your own git account so that the history attributes it to you.
\clearpage
\begin{formbox}
\small
\textbf{Individual Contribution Statement --- Project 2 (PDE BVP)}\\[0.3em]
\textbf{Name:} Han Liu \hfill \textbf{Student ID:} 2364393\\
\textbf{Role this project:} Report integration, conclusions and
reproducibility packaging.

\textbf{1. The specific things I did this project:}
\begin{enumerate}[leftmargin=1.5em,itemsep=1pt,topsep=2pt]
    \item I integrated the five section files into one report: I merged the
    theory (\S1--\S2), the algorithm package (\S3), the experiments (\S4) and
    the conclusions into a single document with one numbering scheme, one
    notation for the face conductances $a_{i+1/2,j}$, and one set of
    cross-references, and I checked that every table and figure is referenced
    from the text and placed after its first mention.
    \item I wrote \S5, which synthesises the three results that the project
    actually establishes --- the conservative flux formulation is what keeps
    the discrete operator symmetric, the harmonic interface conductance is
    exact on the layered benchmark, and the arithmetic mean is symmetric and
    conservatively balanced while its interface flux is still wrong by $36\%$
    --- and states the limitations: one interface geometry, one contrast range,
    and a maximum-norm bound for the convergence theorem that is not sharp in
    two dimensions.
    \item I wrote Appendix~C and the repository packaging: the \path{report/}
    tree with \path{tables/}, \path{figures/} and \path{src/}, the build files,
    and the commit identifiers recorded for provenance. I verified that the
    report builds from a fresh clone with zero errors, zero overfull boxes and
    zero undefined references, so that the submitted PDF is reproducible rather
    than a locally patched artefact.
    \item I ran the final consistency pass over the numbers: every quantity
    quoted in the abstract and in the conclusions was traced back to a table
    generated by the code, duplicated figures between \S3 and \S4 were removed,
    and the captions were rewritten so that each states what the figure is
    evidence for instead of only what it plots.
\end{enumerate}

\textbf{2. The hardest conceptual obstacle I encountered:}
Making a five-author document read as one argument, when the sections arrived
with different conventions: the theory section described the arithmetic mean in
one dimension, the experiment section reported a two-dimensional maximum, and
the abstract quoted one of the two. The obstacle was to see that these were not
inconsistent but described different quantities, and the fix had to be made in
the theory section (by naming which quantity each statement is about), not by
deleting the experiment number.

\textbf{3. One piece of AI-generated output my team used, and what I changed:}
An AI-drafted conclusion claimed that ``both methods show second-order
convergence, which is confirmed on the layered benchmark as well''. That
sentence is wrong twice: the layered benchmark has no asymptotic order to
confirm, because its exact solution is piecewise linear and the harmonic scheme
is nodally exact there, and the arithmetic variant is first order in the nodal
norm. I rewrote the conclusion so that the smooth benchmark carries the order
claim and the layered benchmark carries the exactness claim (Appendix~A,
Audit~1).

\textbf{4. Anything I should have asked for help with earlier:}
I should have asked for a single machine-generated table of all reported
numbers at the start of integration, instead of reconciling numbers quoted in
four different files by hand.

\textbf{Signature:} \underline{\hspace{4.2cm}} \hfill
\textbf{Date:} \underline{\hspace{2.2cm}}
\end{formbox}
